\begin{donnees}
    \begin{itemize}
        \item $m=\qty{100}{\gram}$~; les frottements sont négligeables.
    \end{itemize}
\end{donnees}

On écarte $\mathrm{(S)}$ (Fig.~\ref{fig:solide_ressort}) de sa position
d'équilibre d'une distance $X_{m}$ et on l'abandonne sans vitesse initiale à
l'instant $t_{0}=0$. Le solide $\mathrm{(S)}$ effectue 10 oscillations pendant
la durée ${\Delta t=\qty{3,14}{\second}}$.

\begin{minipage}[t][][t]{0.65\linewidth}
    \begin{enumerate}
        \item Déterminer la valeur de la période propre $T_{0}$.
        \item Déduire la valeur de $K$.
        \item On choisit l'état où le ressort n'est pas déformé comme référence
              de l'énergie potentielle élastique $E_{\mathrm{pe}}$ et le plan
              horizontal contenant $G$ comme état de référence de l'énergie
              potentielle de pesanteur $E_{\mathrm{pp}}$.

              La courbe de la
              Fig.~\ref{fig:2bac_svt_national_2018_normale_phys_ex3_4} représente le
              diagramme d'énergie potentielle élastique $E_{\mathrm{pe}}=f(x)$.
              
              En exploitant le diagramme, déterminer les valeurs de~:
              \begin{enumerate}
                  \item L'amplitude $X_{m}$.
                  \item L'énergie mécanique $E_{m}$ du système oscillant.
                  \item La vitesse maximale du mouvement de ($S$).
              \end{enumerate}
    \end{enumerate}
\end{minipage}
\hfill
\begin{minipage}[t][][b]{0.33\linewidth}
    \begin{figure}[H]
        \centering
        \begin{tikzpicture}
            \def\K{40}
            \begin{axis}[
                    xmin=-0.049,xmax=0.049,
                    ymin=0,ymax=39,
                    width=7cm,height=7cm,
                    xlabel=$x~(\unit{\meter})$,
                    ylabel=$E_{pe}~(\unit{\mJ})$,
                    scaled ticks=false,
                    xtick distance=0.01, ytick distance=4,
                    xticklabels={,,,,,,0,,{0,02}}, yticklabels={,,,8},
                    minor tick num=0,
                ]
                \addplot[domain=-0.04:0.04,samples=200,
                    very thick,blue] {0.5*\K*x^2*1e3}
                    coordinate[pos=0](A) coordinate[pos=1](B);
            \end{axis}
            \foreach \n/\a in {A/200,B/-20}{%
            \draw[very thick] ([yshift=-5mm]\n) -- ([yshift=5mm]\n);
            \foreach \f in {0.02,0.14,...,1}
                \draw[thick]
                    ($([yshift=-5mm]\n)!\f!([yshift=5mm]\n)$) -- ++(\a:.2);
            }
        \end{tikzpicture}
        \caption{}
        \label{fig:2bac_svt_national_2018_normale_phys_ex3_4}
    \end{figure}
\end{minipage}

